\documentclass{situationist-preprint}

\title{Evidence Is Knowledge}
\msid{EPI-21-a}
\msfamily{EPI-21}
\msdiscipline{Epistemology}
\mstype{Principal manuscript}
\mskeywords{evidence; knowledge; E = K; perceptual justification; new evil demon; luminosity}

\begin{document}
\maketitle

\begin{abstract}
What is a subject's evidence? We argue that one's evidence is exactly one's knowledge. Evidence consists of propositions, since it is what one reasons from, what one conditionalizes on and what is consistent or inconsistent with hypotheses. A proposition is part of one's evidence only if one has it in a way that is not accidental, and having a true proposition non-accidentally is knowing it; conversely, whatever one knows is available as a premise and as a condition on probabilities, and so is evidence. The view faces a challenge that its critics press: if evidence is knowledge, the deceived twin of a perceiver has less evidence than the perceiver, and the experiences that ground perceptual knowledge seem left without a role. We answer that experiences justify without being evidence. A perceptual seeming confers prima facie justification on the belief it grounds; when the belief is knowledge, its content becomes evidence. The perceiver and her twin are alike in justification and differ in evidence. Justification has two sources, experiential and evidential, and only the second is knowledge.
\end{abstract}

\section{The problem}

Evidence is what rational belief answers to. A detective's hypothesis is supported by the evidence, a scientist conditionalizes on the evidence, and a juror is told to set aside what is not in evidence. The question what a subject's evidence is has received four main answers. On the \emph{mentalist} view, one's evidence consists of one's current mental states, including experiences and seemings \citep{ConeeFeldman2004}. On a second view, it consists of the propositions one is justified in believing; on a third, of the true propositions one believes on good grounds. On the fourth, defended by \citet{Williamson1997,Williamson2000}, one's evidence is exactly one's knowledge, a thesis abbreviated as E~=~K.

We defend E~=~K. EPI-02 and EPI-20 rely on it: in the good case, a subject's evidence includes that she has hands, or that the animals in the pen are zebras, and so discriminates conclusively between the ordinary hypothesis and the sceptical one. Its critics have pressed one challenge above all \citep{Silins2005,Goldman2009,ComesanaKantin2010}. If evidence is knowledge, perceptual experiences are not evidence, and a subject deceived by a demon, whose experiences are exactly like a perceiver's, has less evidence than the perceiver; yet she seems to be just as rational. We argue that the challenge is answered by distinguishing evidence from justification. Experiences justify, as EPI-19 holds, without being evidence; and the deceived subject is as justified as the perceiver, with less evidence.

\section{The argument}

\begin{principle}[Propositionality]\label{p:prop}
One's evidence consists of propositions.
\end{principle}

Evidence enters into relations that only propositions enter into. A hypothesis is consistent or inconsistent with the evidence; the evidence entails or makes probable a hypothesis; a hypothesis is the best explanation of the evidence; one conditionalizes on the evidence. An experience is not consistent with a hypothesis, and one does not conditionalize on a state; one conditionalizes on the proposition that one is in it, or on a proposition it presents \citep[ch.~9]{Williamson2000}.

\begin{principle}[Non-accidental possession]\label{p:poss}
A proposition is part of one's evidence only if it is true and one has it, as true, in a way that is not accidental.
\end{principle}

A false proposition is not part of anyone's evidence. The detective who relied on a forged letter did not have the evidence he thought he had: what he took as evidence was not so, and his conclusions, however rational, were not supported by the evidence. Nor is every true proposition that one believes with justification part of one's evidence. A subject who believes that it is three o'clock on the basis of a clock that stopped at three, twelve hours earlier, does not have as evidence that it is three o'clock, even though it is; she has as evidence what she knows, that the clock reads three, and her belief about the time is a conclusion from it, which happens to be true. To have $p$ as evidence is to be in a position to treat $p$ as settled, to reason from it outright; and treating a merely luckily true proposition as settled is a mistake even when it does no harm.

\begin{principle}[Availability]\label{p:avail}
Every proposition one knows is available to one as a premise in reasoning and as a condition on probabilities.
\end{principle}

A scientist who knows that the meter reads 4.2 may use that proposition in calculating, may conditionalize on it, and may cite it as evidence for a hypothesis. There is no proposition she knows and is barred from using in these ways, except through considerations of relevance, which concern which evidence bears on a question, not what her evidence is.

\begin{theorem}[E = K]\label{t:ek}
One's evidence is exactly one's knowledge.
\end{theorem}

\begin{proof}
Suppose $p$ is part of one's evidence. By Principles~\ref{p:prop} and~\ref{p:poss}, $p$ is a true proposition that one has, as true, non-accidentally. On the analysis of EPI-01, a belief that is justified, true and safe, so that it could not easily have been false, is knowledge; and a true proposition had non-accidentally is had by a safe and justified belief. So one knows $p$. Conversely, suppose one knows $p$. By Principle~\ref{p:avail}, $p$ is available as a premise and as a condition on probabilities, which is to play the role of evidence; so $p$ is part of one's evidence.
\end{proof}

The argument does not require Williamson's view that knowledge is unanalysable. It is compatible with the analysis of EPI-01, on which knowledge is safe, justified, true belief, and it uses that analysis to identify non-accidental possession of a truth with knowledge.

\section{How experiences justify}

If evidence is knowledge, experiences are not evidence. They are not propositions, and in the bad case, where an experience presents as true what is false, what it presents is not known. The critics infer that E~=~K cannot account for perceptual justification. The inference requires a premise we reject: that all justification is evidential, so that one is justified in believing $p$ only to the extent that one's evidence supports $p$.

Justification has two sources. The first is experiential. If it seems to a subject that $p$, she has prima facie justification for believing $p$ (EPI-19), and an undefeated seeming makes her justified. The seeming does not justify by being a premise from which $p$ is inferred. It justifies by presenting $p$ as true, and the belief it grounds is not inferred from anything. On the structure defended in EPI-05, beliefs grounded in experience are the lowest nodes of the tree of justification, and the experiences are its roots. The second source is evidential: a belief is justified by inference from what one knows, which is one's evidence.

The two sources connect in perception. When a perceptual belief grounded in an experience is safe, it is knowledge, by EPI-01, and its content becomes part of one's evidence, by Theorem~\ref{t:ek}. The experience does not become evidence. It is what puts the subject in a position to know, and so to acquire the evidence. In the good case, then, the subject who sees her hands is justified in believing that she has hands by her experience, knows that she has hands, and has as evidence that she has hands. In the bad case, her deceived twin is justified by her experience to the same degree, since the seemings are the same; she does not know that she has hands, since she has none; and her evidence does not include that she has hands. What the twins share is the seeming and the prima facie justification it confers. What they do not share is the knowledge, and with it the evidence.

This division of labour is what EPI-02 uses against the underdetermination argument for scepticism. The premise that the subject's evidence is the same whether she has hands or is envatted is false: in the good case her evidence includes that she has hands, which favours that hypothesis conclusively, and the envatted subject has less evidence, though she cannot tell that she has less. What is the same is how things seem, which fixes justification, not evidence.

\section{Objections}

\begin{objection}[The new evil demon]
The demon's victim is as rational as we are; her beliefs are as well supported as ours. If her evidence differed from ours, her beliefs would be less well supported. So her evidence is the same as ours, and E~=~K is false \citep{Cohen1984,Silins2005}.
\end{objection}

\begin{reply}
The victim's beliefs are as well justified as ours, and that is the intuition the case elicits. It does not follow that they are as well supported by evidence, unless all justification is evidential. On our view they are justified experientially, by seemings that are the same as ours, and they are less well supported evidentially, since the victim knows less. A defender of E~=~K who identifies justification with evidential support must say that the victim's beliefs are less justified than ours, and at most blameless; we do not follow her there. The intuition of parity concerns justification, which the victim has, not evidence, which she lacks.
\end{reply}

\begin{objection}[Access]
One's evidence must be something one can tell one has; otherwise one cannot be guided by it. Knowledge is not like this: one can know without knowing that one knows, and the victim cannot tell that she lacks the evidence she seems to have \citep{ConeeFeldman2004}.
\end{objection}

\begin{reply}
No non-trivial condition is such that one is always in a position to know whether it obtains, as \citet[ch.~4]{Williamson2000} argues; experiences and seemings are no exception, since there are borderline cases in which one cannot tell whether things seem a certain way. The demand for evidence that is always accessible cannot be met by any candidate, and so it cannot be used to choose between them. What guides the subject is not a second-order verdict about what her evidence is, but the first-order seemings and knowledge that she has, and both are present when she reasons well.
\end{reply}

\begin{objection}[Perceptual knowledge needs evidence]
If perceptual knowledge is to be well founded, it must be based on evidence. But if evidence is knowledge, perceptual knowledge is either based on prior knowledge, which generates a regress, or based on nothing \citep{Brueckner2005}.
\end{objection}

\begin{reply}
Perceptual knowledge is based on experience, which is neither evidence nor nothing. The objection assumes that the only thing a belief can be based on is evidence; we have argued that experiences ground beliefs and confer justification without being evidence. The regress it threatens is the regress of EPI-05, and it stops where it stops there: at perceptual experiences, which are not beliefs and require no further support.
\end{reply}

\section{Conclusion}

One's evidence is one's knowledge. Evidence consists of propositions, is had non-accidentally and is available as a premise; non-accidental possession of a truth is knowledge, and whatever is known is available. Experiences are not evidence, but they justify: a seeming confers prima facie justification on the belief it grounds, and when that belief is knowledge its content becomes evidence. The perceiver and her deceived twin are equally justified and unequally evidenced, and the sceptic who infers sameness of evidence from sameness of experience has confused the two.

\begin{thebibliography}{99}\small
\bibitem[Brueckner(2005)]{Brueckner2005} Brueckner, A. (2005). Knowledge, evidence, and skepticism according to Williamson. \emph{Philosophy and Phenomenological Research}, 70(2), 436--443.
\bibitem[Cohen(1984)]{Cohen1984} Cohen, S. (1984). Justification and truth. \emph{Philosophical Studies}, 46(3), 279--295.
\bibitem[Comesa\~na and Kantin(2010)]{ComesanaKantin2010} Comesa\~na, J. and Kantin, H. (2010). Is evidence knowledge? \emph{Philosophy and Phenomenological Research}, 80(2), 447--454.
\bibitem[Conee and Feldman(2004)]{ConeeFeldman2004} Conee, E. and Feldman, R. (2004). \emph{Evidentialism: Essays in Epistemology}. Oxford University Press.
\bibitem[Goldman(2009)]{Goldman2009} Goldman, A.~I. (2009). Williamson on knowledge and evidence. In P.~Greenough and D.~Pritchard (Eds.), \emph{Williamson on Knowledge}. Oxford University Press.
\bibitem[Silins(2005)]{Silins2005} Silins, N. (2005). Deception and evidence. \emph{Philosophical Perspectives}, 19, 375--404.
\bibitem[Williamson(1997)]{Williamson1997} Williamson, T. (1997). Knowledge and evidence. \emph{Mind}, 106(424), 717--741.
\bibitem[Williamson(2000)]{Williamson2000} Williamson, T. (2000). \emph{Knowledge and Its Limits}. Oxford University Press.
\end{thebibliography}

\end{document}
